% =============================================================================
% Seção 7 — Discussão
% =============================================================================
\section{Discussão}
\label{sec:discussao}

Os resultados da Seção~\ref{sec:resultados} podem ser discutidos sob quatro eixos complementares: a adequação do EGO à formulação atual, o desempenho preditivo do modelo substituto, o comportamento da penalização exterior e as limitações estruturais da formulação implementada.

\paragraph{Adequação do EGO à formulação atual: dois regimes bem delimitados.}
O EGO é classicamente justificado em problemas de otimização tipo caixa-preta com avaliações de alto custo computacional \citepabnt{jones1998efficient}. A função pseudo-objetivo $\Theta$ desta formulação, entretanto, tem baixo custo na implementação vetorizada atual ($\approx \SI{0,1}{\milli\second}$ por avaliação), o que torna a justificativa clássica insuficiente e exige avaliação empírica explícita. O protocolo de dois cenários responde à questão com precisão. Sob \textit{orçamento igual de avaliações reais} (S1), as arquiteturas assistidas por modelo substituto reduzem $\Theta$ frente às buscas não assistidas, e o CBO apresenta as melhores médias nos três casos. Sob \textit{tempo de parede} com avaliação de baixo custo (S2), buscas diretas com orçamento estendido entregam soluções melhores em duas ordens de grandeza menos tempo, pois o custo do EGO/CBO é dominado pelos ajustes sucessivos do GPR e pela maximização interna da aquisição. A auditoria por decomposição (Tabela~\ref{tab:decomposicao}) acrescenta uma terceira leitura: nas instâncias simplificadas atuais, uma linha de base por sapata praticamente atinge ou supera o melhor volume factível global, confirmando que o problema ainda não explora o acoplamento que justificará plenamente a arquitetura assistida. A adoção do EGO nesta etapa deve, portanto, ser lida como investimento metodológico deliberado: a arquitetura é candidata natural para extensões em que cada avaliação se tornará genuinamente custosa — posicionamento conjunto com empacotamento, interação entre bulbos de tensão, recalques e eventuais acoplamentos numéricos —, mas não é uma necessidade computacional da formulação algébrica atual. Essa leitura é coerente com \citeonline{mathern2021multiobjective}, que exploram em projeto estrutural uma assimetria entre objetivos de baixo custo e restrições normativas mais custosas, cenário ainda não plenamente ativado no recorte atual do FundaIA.

\paragraph{Sensibilidade ao fator de penalidade: o que o $R^2$ esconde.}
O estudo controlado de penalidade refina o diagnóstico das etapas preliminares desta pesquisa. A constatação original — de que penalidades elevadas degradam o ajuste do GPR — mostrou-se imprecisa quando aferida por $R^2$ global: sob $\alpha = 10^{6}$, a variância do alvo é dominada pelo próprio termo de penalização e o $R^2$ permanece alto ($\approx 0{,}92$), mascarando o problema. A métrica reveladora é o erro absoluto na \textit{região factível}, onde a função de aquisição efetivamente discrimina candidatos: ali, o RMSE sob $\alpha = 10^{6}$ é cerca de cinco ordens de grandeza superior à escala do volume, contra $\approx 1{,}3\,\si{\meter\cubed}$ sob $\alpha = 10^{1}$. A consequência prática é direta: com penalidade agressiva, a EI passa a ser governada pelo gradiente artificial da penalização, e não pela paisagem do volume. Esse resultado fundamenta quantitativamente a adoção de penalidade moderada na arquitetura penalizada — e motivou a implementação, nesta mesma pesquisa, da alternativa que remove a penalização de todos os alvos de regressão: a \textit{Constrained Bayesian Optimization} da Seção~\ref{sec:met_cbo}, discutida a seguir.

\paragraph{Tratamento explícito de restrições: quando compensa.}
O confronto pareado CBO~$\times$~EGO — idênticos em tudo, exceto no tratamento das restrições — produziu um resultado assimétrico: o CBO melhora $\Theta$ nos três casos, mas perde factibilidade estrita nos casos 1 e 2. A explicação é coerente com a mecânica dos dois métodos: com o ótimo encostado na fronteira de tensão e poucas observações, o GP da restrição suaviza o cruzamento do zero e a probabilidade de factibilidade da Equação~\eqref{eq:pof} pode superestimar a margem — um erro \textit{de modelo}; já a penalização direta com fator finito desloca o ótimo penalizado em relação ao ótimo restrito — um erro \textit{de formulação}. Como há apenas uma instância por dimensionalidade e a restrição de sobreposição está inativa, o resultado não demonstra isoladamente que o ganho decorre da dimensão. Ele indica, de forma mais cautelosa, que o tratamento explícito de restrições merece prioridade na próxima fase, desde que acompanhado por regras de seleção final estritamente factíveis; no regime de posicionamento conjunto ($5N$ variáveis), a restrição de sobreposição se tornará ativa e produzirá acoplamento real entre fundações. Nesse regime, variantes escaláveis com regiões de confiança \citepabnt{eriksson2021scbo} são o caminho declarado quando $N$ crescer.

\paragraph{Escolha do \textit{kernel}: menos decisiva do que se supunha.}
O rastreio controlado de 21 configurações mostrou que 19 delas se concentram em faixa estreita de capacidade preditiva ($R^2 \approx 0{,}93 \pm 0{,}011$), falhando apenas as hipóteses estruturais incompatíveis com a resposta (periódica pura e linear pura). Esse achado é coerente com a interpretação de \citeonline{schulz2018tutorial} — a função de covariância codifica hipóteses de regularidade — mas deve ser interpretado como triagem exploratória, pois apenas três réplicas foram usadas por configuração. Para esta formulação, a escolha fina do \textit{kernel} dentro das famílias suaves usuais parece importar menos do que a escala da penalidade e o tamanho da base amostral. A configuração de produção (Matérn $\nu = 2{,}5$) situa-se na banda superior, e não há evidência prática que justifique substituí-la sem um estudo maior.

\paragraph{Ganho sobre buscas não assistidas e leitura de engenharia.}
Sob o mesmo orçamento de 150 avaliações reais, a redução mediana da melhor arquitetura assistida em relação à busca aleatória uniforme foi de $24{,}7\,\%$, $45{,}7\,\%$ e $71{,}4\,\%$ nos casos 1, 2 e 3, respectivamente; para o EGO penalizado isoladamente, as reduções foram $23{,}4\,\%$, $38{,}8\,\%$ e $63{,}9\,\%$. Esses números substituem, agora com sementes pareadas, orçamento equalizado e teste estatístico pareado, a comparação exploratória com Monte Carlo das etapas preliminares. Cabe a ressalva metodológica já assumida anteriormente: a busca aleatória é um piso de referência, não uma representação da prática profissional, na qual heurísticas de pré-dimensionamento aproximam o ponto de partida de regiões viáveis \citepabnttwo{wang2008economic}{waheed2022practical}. A leitura de engenharia dos arranjos finais é coerente com a formulação adotada: a tensão admissível governa os três casos ($S/R \to 1$), a punção nos contornos $C$ e $C'$ permanece folgada na faixa explorada e, no caso 1, a geometria do pilar alongado ($a_p = \SI{2,10}{\meter}$) ativa a restrição de balanço mínimo na direção $x$.

\paragraph{Penalização linear e factibilidade estrita.}
A aferição com tolerância estrita expôs o compromisso estrutural da penalização exterior com fator finito: soluções finais de todos os métodos carregam violações residuais ($g_k \sim 10^{-3}$--$10^{-1}$), pois o ótimo do problema penalizado com $\alpha = 10$ e $p = 1$ não coincide exatamente com o ótimo restrito. Essas violações não devem ser aceitas como soluções de projeto, nem interpretadas como absorvidas por coeficientes parciais de segurança. Elas importam metodologicamente por três razões: (i) quantificam o efeito previsto pela teoria clássica de penalização; (ii) mostram que a avaliação por $\Theta$ precisa ser acompanhada por métricas de factibilidade e melhor volume estritamente factível; e (iii) estabelecem a linha de base contra a qual esquemas alternativos — penalização adaptativa, regras de factibilidade de Deb, seleção final por dominância de factibilidade e \textit{Constrained Bayesian Optimization} — deverão ser julgados na continuidade da pesquisa.

\paragraph{Limitações da formulação atual.}
Cinco limitações delimitam o alcance das conclusões e motivam os trabalhos futuros da Seção~\ref{sec:conclusoes}:
\begin{enumerate}
    \item \textbf{Posições fixas e quase separabilidade.} As posições $(x_g, y_g)$ são dados de entrada; nos três casos congelados, a restrição de não sobreposição é inativa por construção. Assim, o problema conjunto de $N_{\mathrm{f}}$ sapatas é quase separável em subproblemas tridimensionais. A linha de base por decomposição já quantifica esse efeito e mostra que os melhores volumes factíveis globais estão próximos do ótimo decomposto; logo, os ganhos com o aumento de $N_{\mathrm{f}}$ não podem ser atribuídos causalmente à dimensionalidade.
    \item \textbf{Pré-dimensionamento, não projeto executivo.} A punção é verificada nos dois contornos críticos da NBR~6118 ($C$ e $C'$, este a $2d$ da face, com transferência de momentos), mas a ferramenta não dimensiona armadura de flexão, não verifica cisalhamento unidirecional, não trata ancoragem/detalhamento e não calcula custo total. Essa diferença é material: trabalhos de otimização de sapatas voltados a projeto econômico normalmente incluem armadura, custo e verificações estruturais adicionais \citepabnttwo{nigdeli2018metaheuristic}{waheed2022practical}. A taxa $\rho$ que alimenta $\tau_{Rd1}$ é a mínima normativa, adotada como hipótese de modelagem \citepabnttwo{waheed2022practical}{santos2018punching}.
    \item \textbf{Simplificações geotécnicas e de ações.} A correlação $N_{\mathrm{spt}}$--$\sigma_{\mathrm{adm}}$ permanece como hipótese empírica de pré-dimensionamento. Estudos recentes tratam capacidade de carga e recalque de fundações superficiais por modelos de dados calibrados \citepabnttwo{ahmad2021gpr}{khajehzadeh2022effective}, e \citeonline{fattahi2025settlement} destacam a sensibilidade de previsões de recalque ao SPT; tais trabalhos reforçam que a correlação simples usada aqui não deve ser lida como validação geotécnica final. A formulação de tensões já calcula o peso próprio como função explícita do volume e não aplica os coeficientes legados de $1{,}05$ e $1{,}30$ na comparação com $\sigma_{\mathrm{adm}}$, mas ainda é necessário separar formalmente combinações de serviço e de estado limite último em uma etapa de projeto completo.
    \item \textbf{Independência geotécnica.} Interações por bulbo de tensões, recalques diferenciais e variabilidade geotécnica não são modeladas \citepabnt{juang2013reliability}.
    \item \textbf{Orçamento e generalização.} As conclusões de S1 valem para o regime de 150 avaliações nas três instâncias avaliadas; o caso 3 indica que esse orçamento não esgota o problema, e a extrapolação para dimensionalidades maiores requer novo experimento.
\end{enumerate}

\paragraph{Reprodutibilidade.}
Em alinhamento com a literatura de comparação rigorosa de metaheurísticas \citepabnttwo{gomes2018probabilistic}{morales2020balance}, o protocolo final adotou 30 repetições semeadas por célula experimental, sementes pareadas entre algoritmos, orçamento de avaliações equalizado por cenário, teste não paramétrico pareado com correção para múltiplas comparações e persistência integral de configurações, históricos por avaliação, métricas e metadados de ambiente (versões de bibliotecas e revisão do repositório), com todas as figuras e tabelas regeneráveis por \textit{scripts} determinísticos. A extensão natural do desenho experimental — inclusão de CMA-ES, Differential Evolution como método pareado de orçamento controlado, linhas de base determinísticas e réplicas em problemas de maior porte — permanece como continuidade; em particular, \citeonline{nigdeli2018metaheuristic} reforçam a utilidade de comparar algoritmos não apenas por melhor valor, mas também por média, dispersão e esforço de avaliação em problemas de sapatas.
